\section*{Introduction Générale}
\addcontentsline{toc}{section}{Introduction Générale}
Dans un contexte médical en constante évolution, l'optimisation de la surveillance des patients en salle de réveil constitue un enjeu crucial pour assurer des soins de qualité et réactifs. Actuellement, dans de nombreux établissements de santé en Tunisie, cette surveillance repose encore largement sur des méthodes manuelles et des systèmes centralisés limités, rendant difficile une détection rapide et efficace des situations critiques. Ce constat met en lumière l'importance de moderniser les processus de suivi post-opératoire à l’aide de technologies avancées et interconnectées.

Le présent projet, réalisé au sein du SmartLab de la Faculté de Médecine de Monastir, s’inscrit dans cette démarche d’innovation. Il vise à développer un système de collecte de données en temps réel et de gestion des alertes pour les machines utilisées en salle de réveil. Ce système repose sur une approche combinant l’extraction de données via la reconnaissance optique de caractères (OCR) à l’aide de caméras pour les machines ne disposant pas d’interfaces de communication.

L’objectif principal est d’assurer une prise en charge rapide et ciblée des patients en situation critique, en réduisant les délais d’intervention grâce à une détection automatisée des alertes. En intégrant des technologies web et mobiles, ce système ambitionne d’améliorer significativement la qualité des soins post-opératoires, tout en diminuant la charge de travail manuelle et le risque d’erreurs humaines.

Ainsi, ce projet répond à un besoin concret et urgent dans le domaine de la santé en proposant une solution intelligente, adaptable et évolutive, capable d’apporter une plus grande réactivité et efficacité dans la gestion des soins en salle de réveil.
